\chapter{How Numbers Are Stored}
\label{app:numbers}

\section{The format}
\index{numbers!representation}

A number in \ccprod{} occupies five bytes and is binary floating point --- not
fixed point, not binary-coded decimal, and not a scaled integer.

The format is the one the Commander X16's ROM arithmetic library uses. That is
the whole reason for the choice: an addition costs a jump into ROM bank 4
rather than a software floating-point library that would not fit in the program.

\begin{ccfigure}
\begin{tikzpicture}[x=1pt,y=-1pt,
  by/.style={draw=ccink,line width=0.7pt,minimum height=22pt,inner sep=2pt,
             font=\ttfamily\fontsize{8}{9.5}\selectfont},
  num/.style={font=\sffamily\fontsize{6.8}{8}\selectfont,text=ccink,anchor=south},
  ann/.style={font=\sffamily\fontsize{7.4}{9}\selectfont,text=ccink}]

  \node[by,fill=black!14,minimum width=62pt] (e) at (31,12) {exponent};
  \node[by,fill=black!30,minimum width=16pt] (s) at (70,12) {s};
  \node[by,minimum width=46pt]  (m1) at (101,12) {};
  \node[by,minimum width=62pt]  (m2) at (155,12) {};
  \node[by,minimum width=62pt]  (m3) at (217,12) {};
  \node[by,minimum width=62pt]  (m4) at (279,12) {};

  \node[num] at (31,-1)  {byte 0};
  \node[num] at (93,-1) {byte 1};
  \node[num] at (155,-1) {byte 2};
  \node[num] at (217,-1) {byte 3};
  \node[num] at (279,-1) {byte 4};

  \draw[ccink,line width=0.5pt] (62,26) -- (62,33) -- (310,33) -- (310,26);
  \node[ann,anchor=north] at (186,34) {mantissa, 32 bits, most significant
    first, leading 1 implicit};

  \node[ann,anchor=north west] at (0,52)
    {byte 0 \quad exponent, biased by 128. \textbf{Zero means the value is
     exactly zero.}};
  \node[ann,anchor=north west] at (0,66)
    {bit s \quad\ \ the sign, 1 for negative. It sits where the mantissa's
     always-1 leading bit would be.};
  \node[ann,anchor=north west] at (0,80)
    {value \quad\ \ $\pm\,0.1mmmm\ldots{}_2 \times 2^{\,(\mathrm{exponent}-128)}$};
\end{tikzpicture}
\caption{The five bytes of a number}
\label{fig:mbf}
\end{ccfigure}

Three worked examples:

\begin{cctable}{Value}{Bytes}{1.4in}{3.42in}
1.0 & \lit{81 00 00 00 00} \\
5.0 & \lit{83 20 00 00 00} \\
9.81 & \lit{84 1C F5 C2 8F} \\
\end{cctable}

\section{Range and precision}

\begin{cctable}{Property}{Value}{2.4in}{2.42in}
Significant decimal digits & about 9 \\
Largest magnitude & just under $2^{127}$, about $1.70 \times 10^{38}$ \\
Smallest non-zero magnitude & about $2.94 \times 10^{-39}$ \\
Infinity & does not exist \\
Not-a-number & does not exist \\
Negative zero & does not exist \\
Underflow & silently becomes exactly zero \\
\end{cctable}

Nine digits is the honest figure. Thirty-two mantissa bits are worth about 9.6
decimal digits, and the display and the program's own tests both treat nine as
the number that can be relied on.

\begin{cccaution}
\textbf{Large integers are not exact.} A figure like a share volume of
7,443,220,000 has more significant digits than the format can hold, and lands
on a neighbouring value --- it will come back looking like a transcription
error. The remedy is to scale: store volumes in thousands, money in whole
units, and let a number format put the zeros back.
\end{cccaution}

\section{Binary, not decimal}

\begin{ccnote}
This is binary floating point, so 0.1 plus 0.2 behaves the way it behaves
everywhere else: the result is very slightly not 0.3, and
\fml{=A1+A2=0.3} can be \msg{FALSE} when the cells display \lit{.1} and
\lit{.2}.

Number formats round on display, so a column of money \emph{looks} right; and
\fn{ROUND} in a formula makes the value itself right. Comparing two computed
values for exact equality is the thing to avoid. Compare a difference against a
tolerance instead: \fml{=ABS(A1-A2)<0.001}.
\end{ccnote}

\section{Overflow, and why the checks are strict}

\begin{ccnote}
The ROM's arithmetic library handles overflow and division by zero by jumping
into BASIC's error handler --- and there is no return from that. BASIC would
print its own error, drop to \msg{READY}, and your unsaved workbook would be
gone.

So \ccprod{} range-checks every operation \emph{before} calling the ROM, from
the exponent bytes alone. The checks are deliberately conservative: they may
refuse one value right at the very top of the range that would in fact have
fitted. A spurious \errv{\#NUM!} on an absurd number costs nothing. A crash
costs the workbook.
\end{ccnote}

\begin{cctable}{Operation}{Refused when}{1.35in}{3.47in}
Add, subtract & Either operand has the maximum exponent --- magnitude
$8.5\times10^{37}$ or more --- even if the result would fit \\
Multiply & The exponents sum past the top of the range \\
Divide & The divisor is zero (\errv{\#DIV/0!}), or the quotient would overflow \\
Power & Zero to a negative power (\errv{\#DIV/0!}); a negative base with a
fractional exponent (\errv{\#NUM!}) \\
\fn{MOD} & The divisor is zero \\
\end{cctable}

$x^0$ is 1, including $0^0$.

\section{Rounding}

\begin{cctable}{Operation}{Rounds}{2.0in}{2.82in}
Arithmetic & In the ROM, using a 40-bit accumulator and a guard byte \\
\fn{ROUND}, and every number format & Half away from zero: 2.5 to 3, $-2.5$ to
$-3$. Not banker's rounding \\
\fn{INT} & Down, towards minus infinity: $-2.7$ to $-3$ \\
Conversion to a whole number & Towards zero \\
\end{cctable}

\section{Reading a typed number}

At most \textbf{eleven} significant digits are taken into the mantissa. Further
integer digits still scale the magnitude; further fraction digits are discarded.

A decimal exponent below $-300$ gives \errv{\#NUM!}. The decimal scaling is
applied one factor of ten at a time through the same guarded multiply and
divide, so \lit{1E999} reports \errv{\#NUM!} rather than crashing the machine.

\section{Writing a number out}

In General format the number goes through the ROM's own formatter, with the
leading space that CBM machines reserve for a sign stripped off.

\begin{cctable}{Value}{Displays as}{1.6in}{3.22in}
0 & \lit{0} \\
8.50 & \lit{8.5} --- trailing zeros are not kept \\
1/3 & \lit{.333333333} --- no leading zero \\
$-0.5$ & \lit{-.5} \\
$1.2\times10^{15}$ & \lit{1.2E+15} \\
\end{cctable}

Scientific notation takes over at a thousand million and above, and below one
hundredth. The buffer for a formatted number is 24 bytes.
